%! Setting Up a Test for the Difference of Two Population Means — Unit 7, Lesson 7.8 — Exit Ticket (Answer Key)
\documentclass[10pt]{article}
\usepackage{apstats-article}
\usepackage{apstats-key}   % pulls in -boxes; do NOT also load -boxes
\newcommand{\TallMath}[1]{$\displaystyle #1\rule[-1.4em]{0pt}{3.2em}$}

\begin{document}

\pageheader{Unit 7, Lesson 7.8}{Exit Ticket --- Answer Key}
\namedateperiod

\vspace{0.15in}

A sports scientist believes the mean sprint time for varsity sprinters differs from the
mean sprint time for club sprinters. She randomly selects 18 varsity sprinters
($n_1 = 18$) and 22 club sprinters ($n_2 = 22$) and records each athlete's best
100-meter time. Dotplots of both groups show roughly symmetric distributions with
no outliers.

\begin{enumerate}
  \item Name the appropriate inference method.

        Method: \ans{two-sample $t$-test for $\mu_1 - \mu_2$}

  \item Define $\mu_1$ and $\mu_2$, then state $H_0$ and $H_a$.

        $\mu_1 = $ \ans{mean 100 m time for all varsity sprinters} \qquad
        $\mu_2 = $ \ans{mean 100 m time for all club sprinters}

        $H_0$: \ans{$\mu_1 - \mu_2 = 0$} \qquad $H_a$: \ans{$\mu_1 - \mu_2 \ne 0$}

  \item Verify the conditions for both samples.
        \begin{itemize}
          \item Random: \ansline{Both groups were randomly selected --- Random condition met for each sample.}
          \item Normal ($n_1 = 18 < 30$, $n_2 = 22 < 30$): \ansline{Both $n_1 = 18$ and $n_2 = 22$ are $< 30$; dotplots of both groups show roughly symmetric distributions with no outliers --- Normal condition met for both.}

                Can the test proceed? \ans{Yes} \quad Explain: \ansline{Both samples are random and both sample distributions are roughly symmetric with no outliers, so all conditions are met.}
        \end{itemize}
\end{enumerate}

\begin{teachernote}
Students often forget to check the Normal condition for \emph{both} groups when both $n < 30$.
Reinforce that the condition applies to each sample independently.
\end{teachernote}

\end{document}
